\section{Introduction}
\label{sec:intro}

In modern deep learning, the paradigm of serving specialized models for diverse tasks at the edge has shifted from maintaining multiple independent models to deploying a single, highly parameterized backbone combined with a pool of Task-Specific Parameter-Efficient Fine-Tuning (PEFT) experts, such as LoRA~\cite{hu2021lora} or adapters~\cite{houlsby2019parameter}. At test time, a dynamic router must route input tokens or samples through the network, dynamically blending the expert parameters or activations layer-by-layer based on the incoming query's characteristics~\cite{sable2025, malkin2023}. 

While test-time dynamic model merging offers unprecedented flexibility and memory efficiency for edge-based serving, deploying these systems under real-world sequential streams introduces significant representation noise~\cite{momentummerge}. This noise arises in two orthogonal dimensions of the serving process:
\begin{enumerate}
    \item \textbf{Intra-Sample Depth-Wise Noise:} As a single query propagates through the network, layer-to-layer representation-space variations can cause the router's ensembling coefficients to oscillate wildly. This depth-wise instability is known as \emph{routing jitter}~\cite{sable2025}.
    \item \textbf{Inter-Sample Temporal Noise:} Sequential serving streams (e.g., in a multi-task chatbot or multi-modal edge sensor) are non-i.i.d. and corrupted by high-frequency temporal noise. A router evaluating samples independently (statelessly) will output highly unstable trajectories over consecutive queries, leading to erratic output qualities.
\end{enumerate}

To combat these sources of noise, recent literature has shifted towards \emph{stateful} serving frameworks that propagate routing history across either network depth or sequence time. For instance, Momentum-Merge~\cite{momentummerge} implements a depth-wise Exponential Moving Average (EMA) of ensembling weights, but resets completely for each new sample, remaining vulnerable to temporal stream noise. Symmetrically, PAC-Kinetics~\cite{pac_kinetics} tracks historical state across sequence samples using a learned first-order state-space model optimized via PAC-Bayesian generalization bounds, but enforces a single static ensembling weight across the entire backbone depth, losing the benefits of layer-wise representational alignment. 

The most recent state-of-the-art framework, ChemMerge~\cite{chemmerge}, attempts to solve both problems simultaneously by modeling routing trajectories as non-equilibrium chemical reaction kinetics governed by continuous-time ordinary differential equations (ODEs). While ChemMerge successfully smooths transitions across both dimensions, it introduces significant model parameterization: it relies on biochemical formulations (such as Arrhenius rate equations, activation energies, and virtual-time steps), requires online numerical ODE integration, and relies on five or more sensitive, hand-tuned hyperparameters.

In this work, we adopt the philosophy of \textbf{The Minimalist}. We argue that modern machine learning can sometimes become needlessly complex, and that the best solutions are often those that strip away unnecessary parameters. By deconstructing prior stateful merging frameworks, we reveal that their performance does not stem from biochemical ODEs or complex learned state-spaces. Instead, it is the fundamental mathematical property of \emph{local recursive filtering} that provides noise resistance.

Guided by Occam's razor, we present \textbf{2D-STEM} (2D Spatio-Temporal Exponential Moving Average), a training-free, highly parameter-efficient, bilinear recurrence filter that smooths routing trajectories across both depth and sequence time simultaneously. 2D-STEM avoids the need for numerical ODE solvers, PAC-Bayesian training loops, and biochemical formulations. In their place, we propose a single, elegant bilinear equation that is analytically guaranteed to preserve the probability simplex at serving time, thereby simplifying system deployment and avoiding explicit re-normalization steps. 

However, propagating sequence history over time introduces a fundamental trade-off: a stateful system exhibits \emph{inertial drag} (transition lag), which severely degrades performance when the serving stream abruptly switches tasks (e.g., from MNIST classification to SVHN). To resolve this without introducing the heavy optimization loops of prior work, we propose \textbf{Adaptive Temporal Gating} (ATG). ATG measures on-the-fly stream homogeneity using raw feature coordinates at an early, frozen layer of the backbone. When a task switch is detected, ATG instantly scales down the temporal momentum to zero, letting the system switch tasks immediately and eliminating transition lag.

Our contributions are summarized as follows:
\begin{itemize}
    \item \textbf{Minimalist Deconstruction:} We deconstruct the state-of-the-art stateful serving frameworks, demonstrating that their noise-reduction capabilities can be modeled by a direct, unified 2D bilinear recursive filter, thereby pruning unnecessary biochemical and learning-theoretic complexity.
    \item \textbf{2D-STEM Formulation:} We present a training-free 2D Spatio-Temporal EMA that is analytically simplex-preserving under a simple linear inequality constraint on the momentum coefficients ($\beta_{\text{depth}} + \beta_{\text{temp}, 0} \le 1$).
    \item \textbf{Adaptive Temporal Gating (ATG):} We introduce a zero-overhead, similarity-based gating mechanism that dynamically scales sequence-level momentum to instantly suppress inertial lag under rapid task switches.
    \item \textbf{Controlled Simulation Study:} We implement a rigorous, parameterized 14-layer representation-space simulation environment (the Analytical Coordinate Sandbox, or ACS) to evaluate stateful routers. We show that 2D-STEM reduces excess routing noise above the physical boundary limit by up to $15.8\times$ (and absolute routing jitter by up to $2.75\times$) compared to SABLE in homogeneous streams, while we mathematically analyze and empirically demonstrate the fundamental trade-offs between sequence-smoothing and task-switching latency.
\end{itemize}

The rest of the paper is organized as follows: Section~\ref{sec:related} contextualizes our approach within the PEFT and dynamic merging literature; Section~\ref{sec:method} details the mathematical formulation and simplex-preserving proofs of 2D-STEM; Section~\ref{sec:experiments} presents our extensive empirical evaluation, and Section~\ref{sec:conclusion} concludes with a discussion of our minimalist research philosophy.
